\documentclass[11pt]{article}

\usepackage[T1]{fontenc}
\usepackage{lmodern}
\usepackage[margin=1in]{geometry}
\usepackage{amsmath,amssymb}
\usepackage{booktabs}
\usepackage{graphicx}
\usepackage{xcolor}
\usepackage{microtype}
\usepackage{hyperref}

\hypersetup{
  colorlinks=true,
  linkcolor=blue!45!black,
  citecolor=blue!45!black,
  urlcolor=blue!55!black,
  pdftitle={The gpt5-6 spinor and the dark sector},
  pdfauthor={Reproducible computational synthesis}
}
\graphicspath{{artifacts/figures/}}
\newcommand{\gptfield}{\ensuremath{\mathrm{gpt5\_6}}}

\title{The \texorpdfstring{\gptfield}{gpt5-6} Spinor and the Dark Sector\\
\large Relationships, conclusions, and claim boundaries}
\author{Reproducible computational synthesis}
\date{17 September 2026}

\begin{document}
\maketitle

\begin{abstract}
This note isolates the relationships between the 16-component Dirac field
\gptfield{}, dark energy, and dark matter.  The field is a reducible
four-flavor multiplet whose homogeneous scalar bilinear satisfies
$S=S_0a^{-3}$.  For a nonlinear potential $U(S)$, its stress tensor obeys
$\rho_\Psi=U$, $p_\Psi=SU_{,S}-U$, and
$w_\Psi=-1+d\ln U/d\ln S$.  Consequently, $U=mS$ is dust-like,
$U=\Lambda$ is vacuum-like, and $U=\Lambda+mS$ reproduces a mixed
matter-plus-vacuum background.  The repository's implemented potential instead
reconstructs an evolving CPL dark-energy history while matter remains a
separate fluid.  We state the exact identities, numerical conclusions, and the
perturbative and observational tests still required before this mathematical
connection could become a physical unified-dark-sector model.
\end{abstract}

\tableofcontents

\section{Purpose and conclusion}

This publication separates three kinds of statement: identities derived from
the spinor action, background interpretations produced by choices of the
potential, and physical claims that need calculations beyond the present
repository.  Its central conclusion is:

\begin{quote}
A nonlinear Dirac condensate can reproduce dark-matter-like, vacuum-like, or
evolving-dark-energy-like homogeneous stress tensors through the choice of its
scalar potential.  The 16-component construction supplies four Dirac flavors,
but it does not by itself prove that one field is both dark matter and dark
energy.  In the implemented CPL reconstruction, \gptfield{} is the effective
dark-energy component and matter remains a separate fluid.
\end{quote}

This distinction is essential because an FLRW expansion history measures the
total homogeneous density and pressure.  It does not uniquely identify their
microscopic origin.

\section{Meaning of the 16-component field}

The field is defined by
\begin{equation}
  \Psi=\begin{pmatrix}\psi_1\\ \psi_2\\ \psi_3\\ \psi_4\end{pmatrix}
  \in\mathbb C^{16},
  \qquad
  \Gamma^\mu=I_4\otimes\gamma^\mu,
  \label{eq:multiplet}
\end{equation}
where every $\psi_A$ is an ordinary four-component Dirac spinor.  Thus
\gptfield{} is a reducible four-flavor multiplet, not a new 16-dimensional
irreducible Lorentz representation.  Its flavor-invariant scalar is
\begin{equation}
  S=\bar\Psi\Psi=\sum_{A=1}^{4}\bar\psi_A\psi_A.
  \label{eq:bilinear}
\end{equation}

The current action uses a potential $U(S)$ depending on the total scalar.
Gravity therefore sees this sum at homogeneous level, not a unique dark-sector
role for each flavor.  Sixteen components do not mean sixteen fluids, and
flavor multiplicity alone creates neither negative pressure nor pressureless
matter.

\section{Exact field-to-fluid relationship}

The torsion-free spinor Lagrangian is
\begin{equation}
\mathcal L_\Psi=
\frac{i}{2}\left[
\bar\Psi\Gamma^\mu\nabla_\mu\Psi
-(\nabla_\mu\bar\Psi)\Gamma^\mu\Psi
\right]-U(S).
\label{eq:lagrangian}
\end{equation}
For a homogeneous configuration in flat FLRW spacetime, the Dirac equations
give
\begin{equation}
  \dot S+3HS=0,
  \qquad S=S_0a^{-3}.
  \label{eq:dilution}
\end{equation}
The on-shell energy-momentum tensor is a perfect fluid with
\begin{equation}
  \rho_\Psi=U(S),
  \qquad
  p_\Psi=SU_{,S}-U(S),
  \qquad
  K_D=SU_{,S}=\rho_\Psi+p_\Psi.
  \label{eq:thermodynamics}
\end{equation}
It follows that
\begin{equation}
  w_\Psi=\frac{p_\Psi}{\rho_\Psi}
  =-1+\frac{SU_{,S}}{U}
  =-1+\frac{d\ln U}{d\ln S}.
  \label{eq:slope}
\end{equation}

Equation~\eqref{eq:slope} is the central relationship.  Expansion fixes
$S(a)$, while the logarithmic slope of $U(S)$ fixes the equation of state.  It
also yields
\begin{equation}
  \dot\rho_\Psi+3H(\rho_\Psi+p_\Psi)=0.
\end{equation}
These identities follow from the action and do not depend on a numerical fit.

\section{Potential choices and dark-sector behavior}

\subsection{Power-law classification}

For a power law
\begin{equation}
  U(S)=\lambda S^n,
\end{equation}
Eqs.~\eqref{eq:dilution} and \eqref{eq:slope} give
\begin{equation}
  w_\Psi=n-1,
  \qquad
  \rho_\Psi\propto a^{-3n}=a^{-3(1+w_\Psi)}.
  \label{eq:powerlaw}
\end{equation}
The relevant regimes are:
\begin{itemize}
\item $U=mS$ gives $p=0$, $w=0$, and $\rho\propto a^{-3}$: the
pressureless-matter limit.
\item $U=\Lambda$ gives $p=-\rho$, $w=-1$, and constant density: the
vacuum-energy limit.
\item $U=\lambda S^n$ with $0<n<2/3$ gives
$-1<w<-1/3$: an accelerating quintessence-like background.
\item A positive potential decreasing with $S$ has $SU_{,S}<0$ and
$w<-1$: a phantom-like background.
\end{itemize}

These are homogeneous stress-tensor equivalences.  They do not establish the
perturbation properties, particle interpretation, or stability of the
corresponding spinor theory.

\subsection{Minimal unified background}

The additive potential
\begin{equation}
  U(S)=\Lambda+mS
  \label{eq:additive}
\end{equation}
gives
\begin{equation}
  \rho_\Psi=\Lambda+mS_0a^{-3},
  \qquad
  p_\Psi=-\Lambda,
  \qquad
  w_\Psi(a)=-\frac{\Lambda}{\Lambda+mS_0a^{-3}}.
  \label{eq:mixed}
\end{equation}
It is matter-like at small $a$ and vacuum-like at large $a$.  At homogeneous
level, this single condensate exactly reproduces a dust term plus a
cosmological constant.  More generally, each term in
$U(S)=\sum_n\lambda_nS^n$ contributes the scaling of a constant-
$w=n-1$ fluid.

Equation~\eqref{eq:additive} is the cleanest mathematical connection between
spinor dark matter and spinor dark energy.  It is also degenerate with the
matching standard-fluid background.  Calling it a physical unification
requires its perturbations and interactions to behave correctly; an algebraic
split of $U$ is not evidence for a unified microscopic sector.

\section{Implemented CPL reconstruction}

The repository does not implement Eq.~\eqref{eq:additive} as a unified fit.
Instead, it reconstructs a nonlinear potential from the combined
Unite+BAO+CMB central CPL history \cite{camilleri,chevallier,linder}
\begin{equation}
  w(a)=w_0+w_a(1-a),
  \qquad (w_0,w_a)=(-0.861,-0.60).
  \label{eq:cpl}
\end{equation}
Energy conservation gives
\begin{equation}
  \rho_\Psi(a)=\rho_{\Psi0}
  a^{-3(1+w_0+w_a)}
  \exp\left[3w_a(a-1)\right].
  \label{eq:cpldensity}
\end{equation}
Using $a=(S_0/S)^{1/3}$ produces
\begin{equation}
\begin{split}
  U(S)={}&\rho_{\Psi0}
  \left(\frac{S}{S_0}\right)^{1+w_0+w_a}\\
  &\times\exp\left\{3w_a\left[
  \left(\frac{S_0}{S}\right)^{1/3}-1
  \right]\right\}.
\end{split}
\label{eq:reconstructed}
\end{equation}
The FLRW equation in the calculation is
\begin{equation}
  \frac{H^2}{H_0^2}
  =\Omega_{r0}a^{-4}+\Omega_{m0}a^{-3}
  +\frac{U(S_0a^{-3})}{\rho_{c0}},
  \label{eq:friedmann}
\end{equation}
with $\Omega_{r0}=0.00009$, $\Omega_{m0}=0.305$, and
$\Omega_{\Psi0}=0.69491$.  The matter term is already present.  Therefore,
in the implemented model, \gptfield{} represents effective evolving dark
energy, while dark matter belongs to the separate matter density.  Relabeling
the spinor as all dark matter would double-count matter.  A unified
interpretation is a different model and needs a new fit.

\section{Numerical conclusions}

At the present epoch, the reconstructed field has
\begin{equation}
\begin{aligned}
  \rho_\Psi/\rho_{c0}&=0.69491000,
  &p_\Psi/\rho_{c0}&=-0.59831751,\\
  K_D/\rho_{c0}&=0.09659249,
  &w_\Psi&=-0.86100000.
\end{aligned}
\label{eq:today}
\end{equation}
The central CPL curve crosses $w=-1$ at
\begin{equation}
  a=0.76833333,
  \qquad z=0.30151844.
  \label{eq:crossing}
\end{equation}
Since
$\rho_\Psi+p_\Psi=SU_{,S}=\rho_\Psi(1+w_\Psi)$, this is also the
homogeneous null-energy-condition boundary.  The complete radiation, matter,
and spinor background begins accelerating at
\begin{equation}
  a=0.57194538,
  \qquad z=0.74841869.
\end{equation}
The formal future end of acceleration at $a=1.80120338$ results from extending
the linear-in-$a$ CPL expression outside its fitted domain and is not an
observational forecast.

The numerical implementation independently checks $S=S_0a^{-3}$,
$\rho_\Psi=U(S)$, and Friedmann closure.  The maximum relative closure error is
$1.80\times10^{-12}$; this is a consistency check, not a physical uncertainty.

\begin{figure}[p]
\centering
\includegraphics[width=0.92\textwidth]{gpt5_6_energy_budget.pdf}
\caption{Radiation, separate matter, and the reconstructed spinor density.
The existing matter curve prevents reinterpreting the fitted spinor density as
the same matter without double counting.}
\label{fig:budget}
\end{figure}

\begin{figure}[p]
\centering
\includegraphics[width=0.92\textwidth]{gpt5_6_potential.pdf}
\caption{The reconstructed $U(S)$ and on-shell kinetic density
$K_D=SU_{,S}$.  Their ratio fixes the evolving equation of state through
Eq.~\eqref{eq:slope}.}
\label{fig:potential}
\end{figure}

\clearpage
\section{Conclusions and claim boundaries}

\subsection{Established by the repository}

\begin{enumerate}
\item Four ordinary Dirac flavors form the defined 16-component multiplet and
satisfy the enlarged Clifford algebra.
\item The homogeneous scalar bilinear dilutes exactly as $a^{-3}$.
\item A potential $U(S)$ maps to a perfect-fluid background through
$\rho=U$ and $p=SU_{,S}-U$.
\item Linear, constant, additive, and reconstructed nonlinear potentials
connect the same framework to dust, vacuum energy, mixed dark-sector
backgrounds, and evolving dark energy.
\item The implemented nonlinear potential reproduces the prescribed central
CPL history and closes the Friedmann equation numerically to about $10^{-12}$.
\end{enumerate}

\subsection{Not established}

\begin{enumerate}
\item The expansion history does not identify a spinor uniquely.  The
potential was reconstructed from that history, and other fields or fluids can
share the same $H(a)$.
\item A pressureless homogeneous limit is insufficient for dark matter.  Cold
dark matter must cluster with acceptable effective sound speed and anisotropic
stress and satisfy CMB, lensing, and structure-growth constraints.
\item Negative pressure is insufficient for fundamental dark energy.  The
fermionic perturbations must be ghost-free, gradient-stable, causal, and under
radiative control, especially across the NEC boundary.
\item No flavor-resolved dynamics assigns some of the four spinors to dark
matter and others to dark energy.  Since $U$ depends on total $S$, that
assignment is absent from the present action.
\item No raw cosmological likelihood or model comparison favors \gptfield{}
over standard dark-sector models.
\end{enumerate}

The 16 components provide internal field content and enough freedom to form a
homogeneous condensate; they do not supply an observational signature by
themselves.  The potential controls the background behavior.  The most
defensible result is therefore a verified field-to-background correspondence
and a concrete research hypothesis, not a dark-sector discovery.

\section{Falsifiable route to unification}

A serious unified-dark-sector extension should:
\begin{enumerate}
\item specify whether the potential depends on total $S$ or separate
bilinears $S_A$;
\item define flavor symmetries, masses, cross-couplings, and a dimensionful
condensate normalization;
\item derive the quadratic action and linearized fermionic and metric
perturbation equations;
\item test ghost, gradient, causal, and strong-coupling conditions;
\item calculate effective sound speed, anisotropic stress, and subhorizon
growth;
\item evolve CMB, lensing, and matter-power observables in a Boltzmann solver;
\item fit supernova, BAO, CMB, lensing, and growth data without a separately
double-counted matter term;
\item compare the evidence with Lambda-CDM and simpler evolving-dark-energy
models.
\end{enumerate}

Passing these tests would distinguish a physical unified spinor sector from a
background reparameterization.

\section{Reproducibility}

The full derivation is in \path{docs/gpt5_6_cosmology.md}; the authoritative
implementation is \path{src/gpt5_6_cosmology.py}; the generated values are in
\path{artifacts/gpt5_6_summary.json}; and source lineage is recorded in
\path{PROVENANCE.md}.  Build both reports and run every check from the
repository root with
\begin{verbatim}
bash scripts/build_documentation.sh
bash scripts/verify_all.sh
\end{verbatim}

\begin{thebibliography}{9}
\raggedright

\bibitem{camilleri}
R. Camilleri et al.,
\emph{Supernovae Unite: Combining Pantheon+ and DES-SN5YR},
arXiv:2609.05053v2,
\url{https://arxiv.org/abs/2609.05053}.

\bibitem{caiwang}
Y.-F. Cai and J. Wang,
\emph{Dark Energy Model with Spinor Matter and Its Quintom Scenario},
Class. Quantum Grav. \textbf{25} (2008) 165014,
arXiv:0806.3890,
\url{https://arxiv.org/abs/0806.3890}.

\bibitem{chevallier}
M. Chevallier and D. Polarski,
\emph{Accelerating Universes with Scaling Dark Matter},
Int. J. Mod. Phys. D \textbf{10} (2001) 213--224.

\bibitem{linder}
E. V. Linder,
\emph{Exploring the Expansion History of the Universe},
Phys. Rev. D \textbf{68} (2003) 083504,
arXiv:astro-ph/0208512.

\end{thebibliography}

\end{document}